% SWITCH-CONTACT1
\section*{SWITCH-CONTACT1: passive on/off contacts as a multisign network model}

\textbf{STATUS: COMPUTATIONAL finite experiment plus KNOWN-THEOREM model correspondences. P versus NP remains OPEN.}

\paragraph{Classical model.}
Shannon's relay/contact formulation makes the user's proposed primitive literal: an input
variable controls a make or break contact, and the output is connectivity between two
terminals. Series contacts implement conjunction, parallel contacts implement disjunction,
and normally-closed contacts implement complementation. Thus the contact model does not
abolish Boolean semantics; it relocates them into graph topology.

A two-terminal contact scheme is an undirected multigraph with edges labelled by variables
or their complements, accepting an input iff its enabled-edge subgraph connects the terminals.
Pudl\'ak--\v{Z}\'ak show, in their asymptotic formulation, that polynomially related
contact-scheme size characterizes nonuniform space through
\[
\mathrm{CS}(F)=\mathrm{DCS}(F)=\mathrm{NU\mbox{-}SPACE}(\log F).
\]
Jukna's communication-complexity work gives lower bounds for bounded-depth
contact-gating schemes and makes graph cuts/information flow explicit. These are
model-matched lower bounds, not unrestricted SAT lower bounds.

\paragraph{Multisign augmentation.}
For the exact-weight-eight classifier, let \(H_n\) be the graph whose vertices are signed
three-clauses and whose edges join clauses with overlapping falsifying subcubes. DN8-FRONTIER1
proves
\[
\mathrm{SAT}\iff \exists\{i,j\}\in E(H_n):x_i=x_j=1.
\]
A biclique \(K_{A,B}\subseteq H_n\) has an immediate passive-switch implementation:
place the contacts \(\{x_a:a\in A\}\) in parallel from the source to one junction and
\(\{x_b:b\in B\}\) in parallel from that junction to the sink. Its contact cost is
\(|A|+|B|\), while it realizes every pair in \(A\times B\). Hence the minimum contact
count in this declared depth-two architecture is exactly the minimum weighted biclique-cover
cost of \(H_n\).

This supplies a concrete graph-theoretic meaning for multisign factorization:
\[
\boxed{
\text{signed-clause overlap}
\to
\text{biclique/rectifier factorization}
\to
\text{controlled-switch complexity}.
}
\]

\paragraph{Finite test.}
The reproducible script
\texttt{experiments/universal-dag/switch\_contact\_conflict.py}
enumerates maximal bicliques and applies the standard efficiency lower bound
\[
\mathrm{Cov}(H)\ge |E(H)|/\kappa(H),
\qquad
\kappa(H)=\max_{K_{a,b}\subseteq H}\frac{ab}{a+b}.
\]

For \(D(4,8)\), \(H_4\) has 32 vertices and 96 edges. Its maximal bicliques have
types \(K_{1,6}\), \(K_{2,2}\), and \(K_{6,1}\), so \(\kappa(H_4)=1\).
An exact MILP over all maximal bicliques attains the lower bound:
\[
\boxed{C^{(2)}_{\rm contact}(4,8)=96}.
\]
The optimum consists of 24 \(K_{2,2}\) bicliques and covers every conflict edge exactly once.

For \(D(5,8)\), \(H_5\) has 80 vertices and 960 edges. Exhaustive maximal-biclique
enumeration gives
\[
\kappa(H_5)=\frac{42}{13},
\]
attained by \(K_{6,7}\)/\(K_{7,6}\). Therefore
\[
C^{(2)}_{\rm contact}(5,8)\ge
\left\lceil\frac{960}{42/13}\right\rceil
=298.
\]
A deterministic greedy weighted cover gives a 605-contact construction, hence
\[
\boxed{298\le C^{(2)}_{\rm contact}(5,8)\le605}.
\]
The upper bound is not claimed optimal.

\paragraph{Critical accounting correction.}
Switch count alone is not a safe resource. If static wires, junctions, rectifiers, or repeated
control fanout are free, a complicated relation can be hidden in the fixed topology.
Accordingly the proposed network resource should be vector-valued:
\[
\mathcal R_{\rm switch}
=(C,R,L,D,B,T,S),
\]
where \(C\) is controlled-contact count, \(R\) routing/topology cost, \(L\) maximum
copies/load of one input control, \(D\) conductive path depth, \(B\) cut breadth,
and \(T,S\) are settling time and retained relay state when intermediate relays are allowed.

This is the main augmentation over a Boolean-gate count. It exposes locality, fanout, routing,
and state rather than charging only abstract operations.

\paragraph{Depth hierarchy and hostile control.}
Depth-two biclique covers are not general contact networks. Rectifier-network literature
shows that weighted biclique-cover cost can exceed unrestricted rectifier-network size by
a polynomial power. Thus any lower bound that survives only depth two is analogous to an
OBDD-only result and must be attacked by deeper routing.

Likewise, allowing intermediate relay outputs to control later contacts introduces sequential
state and returns toward general circuit/finite-state computation. Such power must be charged
rather than silently excluded.

Fixed \(m=8\) remains a deliberate easy control: its conflict graph yields polynomial
representations. The next meaningful experiment is to compute the deeper rectifier/contact
compression of \(H_4,H_5\), then test whether multisign rank, communication cuts, or
breadth predict the minimum topology after biclique factoring is no longer sufficient.

\begin{thebibliography}{9}
\bibitem{Shannon38}
C.~E.~Shannon,
\emph{A Symbolic Analysis of Relay and Switching Circuits},
Trans. AIEE 57 (1938), 713--723.

\bibitem{Jukna87}
S.~P.~Jukna,
\emph{Lower Bounds on Communication Complexity},
Mathematical Logics and Its Applications 5 (1987), 22--31.

\bibitem{Jukna88}
S.~P.~Jukna,
\emph{Entropy of Contact Circuits and Lower Bounds on Their Complexity},
Theoretical Computer Science 57 (1988), 113--129.

\bibitem{Potechin10}
A.~Potechin,
\emph{Bounds on Monotone Switching Networks for Directed Connectivity},
ECCC TR09-142, rev. 2010.

\bibitem{IvanEtAl14}
S.~Iv\'an, A.~D.~Lelkes, J.~Nagy-Gy\"orgy, B.~Sz\"or\'enyi, and G.~Tur\'an,
\emph{Biclique Coverings, Rectifier Networks and the Cost of $\varepsilon$-Removal},
DCFS 2014, LNCS 8614, 174--185.

\bibitem{PudlakZak83}
P.~Pudl\'ak and S.~\v{Z}\'ak,
\emph{Space Complexity of Computations},
preprint, Charles University, 1983.
\end{thebibliography}

% END_OF_NOTE::SWITCH_CONTACT1
